Ett \emph{teleobjektiv} består av en konvex lins med fokallängden $f_1 = 20\,\mathrm{cm}$ följd av en konkav lins med fokallängden $f_2 = -10\,\mathrm{cm}$. Avståndet mellan linserna är $d = 15\,\mathrm{cm}$.
\begin{enumerate}
  \item Beräkna systemets ekvivalenta fokallängd $f_e$.
  \item Objektivet riktas mot ett objekt oändligt långt bort. Hur långt bakom den konkava linsen ska kamerasensorn placeras för att bilden ska bli skarp?
  \item Hur långt blir objektivet totalt, från första linsen till sensorn? Jämför med en enkel lins som har samma fokallängd som $f_e$.
  \item Månen upptar en vinkel på ungefär $0{,}52^\circ$ på himlen. Hur stor blir månens bild (i diameter) på sensorn med teleobjektivet? Förklara varför konstruktionen är användbar.
\end{enumerate}
